\printedpage{44}
% head: A. Grothendieck : Dérivateurs, Ch. I, Édition des Universités de Montpellier II et Paris VII
the passage from $X$ to $X^\circ$ amounts to interchanging $p_X$ and $q_X$. The fiber of $\Phi(X)$ at $x \in \Ob X$ is $\simeq (X/x)^\circ$, that at $y^\circ \in \Ob X^\circ$ is $\simeq y\bs X$
\[
\left\{
\begin{aligned}
&p_X^{-1}(\{x\}) \simeq (X/x)^\circ\\
&q_X^{-1}(\{y^\circ\}) \simeq y\bs X
\end{aligned}
\right.
\]

N.B. : $p_X$ is $W$-smooth \texttt{[rather $W$-proper]} with $W$-aspheric fibers, hence in $W^u$ \texttt{[$W^u =$ set of morphisms which are in $W$ and remain so after any base change]}, $q_X$ is $W$-proper with $W$-aspheric fibers, hence in $W^u$.

\msp{page 74}

One checks
\[
\left\{
\begin{aligned}
&\Phi(X)/x \simeq \Phi(X/x) &\qquad& x \in \Ob X\\
&\Phi(X)/x^\circ \simeq \Phi(x\bs X) &\qquad& x^\circ \in \Ob X^\circ
\end{aligned}
\right.
\]

I say that $p_X$, $q_X$ are $\D$-\textit{strict equivalences} and \textit{costrict.} By the preceding formulas, this means that the $\Phi(X/x)$, $\Phi(x\bs X)$ are $\D$-contractible. Thus that if $Z$ has a final object, $\Phi(Z)$ is $\D$-contractible. But let us consider
\[
\begin{tikzcd}
\Phi(Z) \arrow[dr, "q_Z"] & \\
& Z^\circ
\end{tikzcd}
\quad ;
\]
it suffices to prove that this is a strict $\D$-equivalence (since $Z^\circ$, having an initial object, is $\D$-contractible), hence that the $z_0\bs\Phi(Z)$ are $\D$-contractible. We have by the preceding formulas % [sic] contractil -> contractile, adjectif ailleurs écrit « contractile »
\[
z_0\bs\Phi(Z) \simeq \Phi(z_0\bs Z) = \Phi(T) \quad ,
\]
where this time $T$ has an initial object $\emptyset$ and a final object $e$. Let us then consider the object
\[
\emptyset \longrightarrow e
\]
of $\Phi(T)$, it is a final object, since for every $x \to y$ in $\Phi(T)$, there is a unique arrow
\[
\begin{tikzcd}
x \arrow[r] & y \arrow[d] \\
\emptyset \arrow[u] \arrow[r] & e
\end{tikzcd}
\]
% foot: Version 8 juillet 2001

\printedpage{45}
% head: A. Grothendieck : Dérivateurs, Ch. I, Édition des Universités de Montpellier II et Paris VII
\msp{page 75}

To summarize

\textbf{Proposition 1 :} \textit{$p_X$ and $q_X$ are strict} \textsf{hot}\textit{-equivalences,} i.e. \textsf{Hot}\textit{-aspheric (hence also strict $\D$-equivalences if $\Phi(X)$ is in the domain of $\D$) with fibers $(X/x)^\circ$ and $y\bs X$ respectively (hence with contractible fibers).}

\textbf{Corollary 1.} \textit{For $X$ to be $\D$-contractible, it is necessary and sufficient that $X^\circ$ be so.}

\hspace*{\parindent}More generally :

\textbf{Corollary 2.} \textit{Let $f : X \to Y$. For $f$ to be a $\D$-equivalence, it is necessary and sufficient that $f^\circ$ be so.}

\hspace*{\parindent}We shall use that the diagram
\[
\begin{tikzcd}
& \Phi(X) \arrow[dl, "p_X"'] \arrow[dr, "q_X"] & \\
X & & X^\circ
\end{tikzcd}
\]
is functorial in $X$. We therefore have
\[
\begin{tikzcd}
& \Phi(X) \arrow[rr, "\Phi(f)"] \arrow[dl, "p_X"'] & & \Phi(Y) \arrow[dl, "p_Y"] \arrow[dd, "q_Y"] \\
X \arrow[rr, "f" near end] \arrow[from=1-2, to=3-2, "q_X"' near end, crossing over] & & Y & \\
& X^\circ \arrow[rr, "f^\circ"'] & & Y^\circ
\end{tikzcd}
\]
the two squares commutative. Since $p_X$ and $p_Y$ are $\D$-equivalences, $f$ is a $\D$-equivalence if and only if $\Phi(f)$ is, and likewise for $f^\circ$. These arguments assume that if $X$ is in the domain of $\D$, $\Phi(X)$ is too. This is the case if \Dia{} consists of finite ordered categories.
% foot: Version 8 juillet 2001

\printedpage{46}
% head: A. Grothendieck : Dérivateurs, Ch. I, Édition des Universités de Montpellier II et Paris VII
\msp{page 76}

\subsection*{$\D$-fibrations}

\textbf{Proposition 1.} \textit{Let $f : X \to Y$ be in} \Dia\textit{, equivalent conditions :}

\hspace*{\parindent}a) \textit{For every $y \in Y$, the inclusion $X_y \to X/y$ is a $\D$-equivalence, and for every arrow $y_1 \to y_0$ of $Y$, $X/y_1 \to X/y_0$ is a $\D$-equivalence.}

\hspace*{\parindent}b) \textit{For every} constant \textit{coefficient $F$ on $Y$, $f_*f^*F$ is locally constant,}\footnote[33]{By this we mean that for every $u : y \to y'$ in $Y$, $G_{y'} \to G_y$ is \texttt{[an]} isomorphism. Without additional axioms on $\D$, it is not clear that this implies that the ``restrictions'' of $G$ to the $Y/y$ or $y\bs Y$, are isomorphic to constant coefficients. See however axiom \textbf{Der 3} a$'$), page 97, which ensures that it comes to the same thing.} \textit{and ``is computed fiber by fiber''} i.e. \textit{the canonical homomorphism}
\[
f_*f^*(F)_y \overset{\sim}{\longrightarrow} H^*_{\D}(X_y, F|X_y)
\]
\texttt{[is an isomorphism]}.

These conditions are equivalent, and imply the conditions of (b) for every $F$ \textit{locally} constant on $Y$.

The first condition in (a) means that for every $y \in Y$, and every $M \in \Ob \mathcal{A}$,
\[
H^*_{\D}(X/y, M) \longrightarrow H^*_{\D}(X_y, M)
\]
is an isomorphism. Now this arrow is interpreted as
\[
f_*(M_X)_y \longrightarrow H^*_{\D}(X_y, M)
\]
is an isomorphism, \textit{i.e.} that $f_*(M_X)$ is computed fiber by fiber. Since $M_X \simeq f^*(M_Y)$, this is the second condition (b). The second condition in (a), for $y_1 \to y_0$ means that for $M$ in $\mathcal{A}$,
\[
H^*_{\D}(X/y_0, M) \longrightarrow H^*_{\D}(X/y_1, M)
\]
is \texttt{[an]} isomorphism, now this arrow is interpreted as
\[
f_*(M_X)_{y_0} \longrightarrow f_*(M_X)_{y_1} \quad ,
\]
that this is always an isomorphism, this is the first condition (b). That (b) also holds for $F$ locally constant on $Y$ is immediate, since this is a condition of a local nature with respect to $Y$.

\msp{page 77}

The dual proposition calls on the functors
\[
y_0\bs X \longrightarrow y_1\bs X \qquad \text{and} \qquad X_y \longrightarrow y\bs X \quad ,
\]
and the sheaves $f_!f^*M$, for $M$ constant, or locally constant on $Y$.
% foot: Version 8 juillet 2001

\printedpage{47}
% head: A. Grothendieck : Dérivateurs, Ch. I, Édition des Universités de Montpellier II et Paris VII
\textbf{Definition 1.} If the equivalent conditions of proposition 1 are satisfied, we say that $f$ is a $\D$-\textit{cohomological fibration} (or $\D$-\textit{left fibration}). If the dual conditions are satisfied, we say that it is a $\D$-\textit{homological fibration} (or $\D$-\textit{right fibration}).

N.B. Cohomological $\D$-fibrations are stable under any base change which is a local isomorphism. Homological $\D$-fibrations under any base change which is a local co-isomorphism. But these are not notions stable under \textit{every} base change\footnote[34]{? give example!}. But the notion of a cohomological $\D$-fibration is stable under $\D$-smooth base change, and that of a homological $\D$-fibration under $\D$-proper base change.

\hspace*{\parindent}I am going to try to construct an example where base change gets botched, by setting $Y = \boldsymbol{\Delta}_1$. Giving $X$ over $\boldsymbol{\Delta}_1$ amounts to giving an open subcategory $X_0$, and its closed complement $X_1$
\[
\begin{tikzcd}
& X & \\
X_1 \arrow[ur, hook] \arrow[d, dash, dashed] & & X_0 \arrow[ul, hook'] \arrow[d, dash, dashed] \\
1 & & 0 \arrow[ll]
\end{tikzcd}
\quad .
\]
To say that $X \to Y = \boldsymbol{\Delta}_1$ is a cohomological $\D$-fibration simply means
\msp{page 78}
that the restrictions
\[
X_0 \overset{j}{\longrightarrow} X \quad , \qquad X_1 \overset{i}{\longrightarrow} X
\]
are $\D$-equivalences. Let us take a base change $Y' \to Y = \boldsymbol{\Delta}_1$, \textit{i.e.} a $Y'$ with $Y'_0, Y'_1$
\[
\begin{tikzcd}
& Y' & \\
Y'_1 \arrow[ur, hook] \arrow[d, dash, dashed] & & Y'_0 \arrow[ul, hook'] \arrow[d, dash, dashed] \\
1 & & 0 \arrow[ll]
\end{tikzcd}
\quad .
\]
We see that the condition $Y'$ smooth over $\boldsymbol{\Delta}_1$ means that for every $y'_1 \in \Ob Y'_1$, the category $Y'_0/y'_1$ of objects $y_0$ of $Y'_0$ over $y'_1$ (\textit{i.e.} equipped with $y_0 \to y'_1$) is contractible. We see that if $Y'_1 \neq \emptyset$ then $Y'_0 \neq \emptyset$ -- but an example with $Y'_0 = \emptyset$, where $Y'_1 = \emptyset$, is not interesting
% foot: Version 8 juillet 2001
\printedpage{48}
% head: A. Grothendieck : Dérivateurs, Ch. I, Édition des Universités de Montpellier II et Paris VII
(for a projection $I \times J \to I$ is always a cohomological $\D$-fibration). Let us take for example for $Y'$
\[
\begin{tikzcd}
& a \arrow[dl] \\
b & \\
& a' \arrow[ul] \\
\underbrace{\qquad}_{\displaystyle Y'_1} & \underbrace{\qquad}_{\displaystyle Y'_0}
\end{tikzcd}
\quad ,
\]
an amalgamated sum of two copies of $Y = \boldsymbol{\Delta}_1$, hence $X' = X \times_Y Y'$ can be viewed as an amalgamated sum $X \amalg_{X_1} X$. But that seems \textit{not} to be a counterexample. We would need $X_1 \to X'$ \textit{not} to be a $\D$-equivalence.

\msp{page 79}

\textbf{Proposition 2.} \textit{Let $f : X \to Y$ be in} \Dia. \textit{Equivalent conditions :}

\hspace*{\parindent}a) \textit{For every $Y'' \overset{p}{\to} Y' \to Y$, whence}
\[
\begin{tikzcd}
X'' \arrow[r, "q"] \arrow[d] \arrow[dr, phantom, "\text{cart}"] & X' \arrow[r] \arrow[d] \arrow[dr, phantom, "\text{cart}"] & X \arrow[d, "f"] \\
Y'' \arrow[r, "p"'] & Y' \arrow[r] & Y
\end{tikzcd}
\quad ,
\]
\textit{if $p$ is a $\D$-equivalence,} \texttt{[then]} \textit{so is $q$.}\footnote[35]{One must assume $Y', Y''$ \textsf{Hot}-aspheric, cf. page 96, and assume axiom Der 3 $a'$) of page 97 to be satisfied.}

\hspace*{\parindent}b) \textit{$f$ is a ``universal'' cohomological $\D$-fibration,} i.e. \textit{is a cohomological fibration, and remains so after any base change.}

\hspace*{\parindent}b$'$) \textit{Condition dual to} (b) \textit{: $f$ is a universal homological $\D$-fibration.}\footnote[36]{b) $\Rightarrow$ a) \underline{botched}, \textit{cf.} marginal annotation,\footnotemark[37] \textit{cf.} counterexample page 89. But it works for the derivators $\D_i$ ($i \in \mathbb{N}$) associated with the ``basic localizers'' $W_i$, \textit{i.e.} with the theory of homotopy $i$-types.}\footnotetext[37]{The implication b) $\Rightarrow$ a) not proved, only in a weaker form :\par \hspace*{\parindent}a$'$) if $Y''$, $Y'$ are $\D$-aspheric and 1-connected, then $X'' \to X'$ is a $\D$-equivalence. (It even suffices to assume it for $Y'$, $Y''$ having final objects (to imply b) or having initial objects (to imply b$'$).)}

\textit{a}) $\Rightarrow$ \textit{b}) : Let us assume \textit{a}), let $Y' \to Y$ be a base change, whence
\[
\begin{tikzcd}
X' \arrow[r] \arrow[d, "f'"'] & X \arrow[d, "f"] \\
Y' \arrow[r] & Y
\end{tikzcd}
\quad .
\]
% foot: Version 8 juillet 2001